%%==================================================
%% conclusion.tex for BIT Master Thesis
%% modified by yang yating
%% version: 0.1
%% last update: Dec 25th, 2016
%%==================================================
\chapter{纳米尖端电化学刻蚀制备集成系统设计与实现}
\enchapter{Design and Implementation of Integrated System for Nano-tip Electrochemical Etching}
\label{chap:conlution}

\section{总结}
\ensection{Conclusions}
深度学习的发展为各个研究领域注入了新活力，在计算蛋白领域，基于深度学习的结构预测与生成取得了显著进展，特别是AlphaFold2的提出，使得通过计算手段预测蛋白结构的准确率大大提升，进而促进了蛋白功能预测与设计的发展。然而在实际应用中，现有的深度学习方法仍然面临一些挑战：现有方法难以有效融合蛋白质中序列和结构的多模态信息，且由于具备湿实验标签数据集规模较小的局限性，难以对特定生物功能例如亲和力进行准确预测，进而在蛋白设计中也难以高效得到高亲和力的结果。而随着大模型的发展和相关技术的成熟，通过预训练获取领域数据的潜在特征再迁移到下游任务中能够有效解决特定数据集规模小的问题。且图Transformer的发展也带来了多模态融合的新思路。因此，本文基于以上挑战和现有条件，开展了以下内容的研究：

（1）本文针对现有亲和力预测方法在数据规模小、特征提取具有片面性以及多模态信息融合不足等问题，提出了一种基于大规模预训练特征和图Transformer架构的PPI-Graphormer模型。具体的实现包括：使用ESM-2蛋白语言模型提取序列层面的进化保守性特征，结合ESM-IF1模型生成的几何感知的结构特征，全面捕捉了蛋白质的潜在高维特征。然后基于分子间作用力构建了氨基酸之间的边关系，并基于此构建了三种注意力的偏置项：氨基酸对类型编码（AA-Type Encoding）：编码210种氨基酸对类型的匹配关系；相互作用力编码（Interact Encoding）：统计氢键、卤键等作用类型；界面掩码（Interface Mask）：通过残基间距离阈值（<7Å）定义界面区域，增强对关键区域的关注，从而搭建了一个图Transformer模块——PPI-Graphomer。最后使用跳跃连接联合图Transformer输出和原特征，使用MLP进行亲和力回归预测。本项研究在PDBbind-PP数据集上完成训练，并在广泛使用的测试集上进行验证。实验结果表明，本文提出的基于PPI-Graphomer的方法在综合数据集上取得了PCC=0.633，MAE=1.57的结果，优于现有基于机器学习和单纯使用AlphaFold2置信度指标的方法。此项研究有助于我们更加深入理解生命活动机制，同时可以运用于潜在的药物研发，有利于药物的高通量筛选。本项工作已被《BMC Bioinformatics》接收，文章名为《PPI-Graphomer: Enhanced Protein-Protein Affinity Prediction Using Pretrained and Graph Transformer Models》。

（2）本文进一步将搭建的亲和力预测模型PPI-Graphomer运用到目前优秀的结合蛋白生成平台BindCraft上，包括以下两个方面工作：首先利用强化学习对ProteinMPNN进行微调，将PPI-Graphomer作为环境（Environment），ProteinMPNN作为智能体（Agent），其生成序列过程作为动作（Action），生成序列与数据中原有结构信息输入到PPI-Graphomer中，产生的打分作为奖励（Reward）反馈给ProteinMPNN，并使用策略梯度进行参数更新；同时根据已有数据中PPI-Graphomer打分与Rosetta打分高度相关的现状，使用PPI-Graphomer打分替换Rosetta能量打分，从而提升平台设计结合蛋白的效率。本文通过Rosetta界面分析和FoldX打分分别检验改进前后的亲和力和稳定性。实验结果表明，微调后的ProteinMPNN产生的序列在PPI-Graphomer打分和Rosetta打分上均优于改进前的模型，说明该训练有效，能够生成较高亲和力的序列。同时改进后的序列在FoldX打分上没有明显下降，说明其生成结果的稳定性能基本保持不变。而当使用PPI-Graphomer替换BindCraft中的Rosetta能量打分后，有效缩减了运行时间，且设计的结合蛋白通过Rosetta界面分析打分基本不变，保持了原有的亲和力水平。



\section{展望}
\ensection{Prospects}
本研究通过大规模预训练模型和图Transformer架构实现蛋白-蛋白亲和力的准确预测，并将其运用在BindCraft中以提升结合蛋白生成性能，然而也存在一些可以改进的地方，主要包括以下几点：

1. 在蛋白质亲和力预测中，当前模型训练依赖PDBbind等实验解析的复合物数据集，其规模相较于其他深度学习任务远远不足，且实验解析的pdb复合物结构往往偏向稳定复合物，而较少关注一些无序区形成的柔性结构，同时对瞬态结合等弱相互作用的捕捉能力较弱。而通过利用主动学习等手段，可以在小批量湿实验的帮助下，搭建开发预测-实验验证这样的闭环系统,形成数据生成-模型优化-实验验证的迭代增强工作流，从而进一步增强对难度较高蛋白的预测能力。

2.蛋白质相互作用研究涉及亲和力预测、结合位点预测、结合蛋白生成等多项任务，而这些任务的关键点都是对相互作用进行建模。未来可构建多任务学习框架，利用不同下游任务的不同标签训练模型，从而实现通过统一的特征提取完成不同领域间的知识融合和迁移，同时也可以在预训练模型基础上进行小样本微调，快速适应不同场景下的任务需求。


